\textbf{COLUNI 2018 -- Questão 10}\par\medskip A partir da leitura do livro de Carolina Maria de Jesus e dos Contos de Conceição Evaristo, analise as afirmativas abaixo:

\par

I. As autoras retratam a realidade dos marginalizados no Brasil, Carolina Maria de Jesus enfatiza a situação de vulnerabilidade social nas favelas, e Conceição Evaristo as condições de pobreza e violência urbana que atingem os afro-brasileiros.

II. Os textos são escritos a partir do lugar de fala da mulher negra que é também o lugar de experiência e de vivência das autoras, que se baseiam em suas histórias de vida para escrever.

III. As autoras apresentam, em seus textos, situações e personagens que evidenciam a realidade sofrida dos afro-brasileiros, e contar essas histórias foi um meio encontrado por elas para superar as dificuldades econômicas.

\par

Está CORRETO o que se afirma em:\begin{enumerate}
\item I e II, apenas.
\item II, apenas.
\item II e III, apenas.
\item I, II e III.
\end{enumerate}
